\chapter{一元函数微分学}\label{ch:2}
\begin{tcolorbox}[title=章首知识地图,colframe=Ink]
导数定义 $\to$ 求导规则 $\to$ 高阶结构 $\to$ 中值定理 $\to$ Taylor近似 $\to$ 单调、极值、凹凸与几何。
\end{tcolorbox}
\begin{examtip}考研主线：复合、隐式、参数求导，高阶导数，中值定理，单调极值与凹凸。常见题型：构造辅助函数证明等式、不等式，分析完整函数图像。竞赛拓展：插值余项、二阶包络与导数估计。上述为学习任务分层，具体范围以所报考科目的当年大纲为准。\end{examtip}
\section{导数、微分与复合求导}\label{sec:2:1}
\subsection{概念、定理与推导}

\begin{definition}{可导与可微}设$f$在$a$的一个开邻域有定义。若差商有有限实极限$f\prime(a)=\lim_{h\to0}\frac{f(a+h)-f(a)}{h}$，称$f$在$a$可导；无穷极限不称为通常意义的导数。区间端点可另定义单侧导数。若$\Delta f=A\Delta x+o(\Delta x)$，称在$a$可微，微分为$\dd f=A\dd x$。一元情形可微等价于可导，并蕴含连续。\end{definition}
\begin{theorem}{求导规则}设所涉及函数在对应点可导，则$ (uv)\prime=u\prime v+uv\prime$，$(\frac{u}{v})\prime=\frac{u\prime v-uv\prime}{v^2}$（$v\ne0$）；复合导数$(f\circ g)\prime=f\prime(g)g\prime$。若$f$在$a$附近的区间上连续且严格单调，$f\prime(a)\ne0$，则反函数在$b=f(a)$处可导，且$(f^{-1})\prime(b)=\frac{1}{f\prime(a)}$。\end{theorem}
\begin{proof}乘积增量拆成$u\Delta v+v\Delta u+\Delta u\Delta v$；最后一项除以$\Delta x$趋零。复合微分将外层增量写成$f\prime(g)\Delta g+o(\Delta g)$即可。\end{proof}
若$F$在$(a,b)$附近属于$C^1$，$F(a,b)=0$且$F_y(a,b)\ne0$，则该点附近存在唯一经过$(a,b)$的$C^1$隐函数$y=y(x)$，并有$y\prime=-\frac{F_x}{F_y}$。参数方程$x=x(t),y=y(t)$在点附近属于$C^1$且$x\prime(t)\ne0$时，$\frac{\dd y}{\dd x}=\frac{y\prime(t)}{x\prime(t)}$。正函数幂式可先取对数。

\begin{method}{导数、微分与复合求导}条件：分式分母、对数真数与反函数导数均须合法。步骤：画复合层次，逐层求导；隐式关系两边求导。易错点：将参数导数与关于$x$的导数混同。\end{method}
\subsection{乘法求导公式全表：每次只对一个因子求导}\label{r2:product}
以下 $u,v,w$ 均为 $x$ 的函数，$c$ 为常数。使用一阶公式只需各因子在所求点可导；高阶公式则要求相应高阶导数存在。乘积公式本身不要求各因子非零。

\begin{theorem}{两因子与多因子的一阶乘法求导}
\begin{align*}
 (cu)'&=cu',\\
 (uv)'&=u'v+uv',\\
 (uvw)'&=u'vw+uv'w+uvw',\\
 \left(\prod_{j=1}^m u_j\right)'&=
 \sum_{k=1}^m\left(u_k'\prod_{\substack{j=1\\j\ne k}}^m u_j\right).
\end{align*}
多因子每一项只对一个因子求导，其余因子保持原样；共 $m$ 个因子就有 $m$ 个加项，不能把全部因子的导数相乘。
\end{theorem}
\begin{proof}
两因子的增量可写为 $\Delta(uv)=v\Delta u+u\Delta v+\Delta u\Delta v$。除以 $h=\Delta x$ 后，前两项趋于 $u'v+uv'$；最后一项为 $(\Delta u/h)\Delta v$，由于可导蕴含连续，极限为零。多因子公式从两因子公式归纳得到。
\end{proof}

\begin{warning}
一般 $(uv)'\ne u'v'$，也不是只求 $u'v$ 或只求 $uv'$。例如 $u=v=x$，正确导数为 $2x$，而 $u'v'=1$。在 $x=0$ 时两个因子均为零，乘积公式仍成立；需要非零条件的是后面的对数写法，而不是乘积法则。
\end{warning}

\begin{example}{两因子乘积：先分组，再保留两项}{基础训练}
对 $y=(x^2+1)\sin x$，取 $u=x^2+1,v=\sin x$，有
\[
 y'=2x\sin x+(x^2+1)\cos x.
\]
再求一次导数时，两项都需使用乘积法则：
\[
 y''=2\sin x+4x\cos x-(x^2+1)\sin x.
\]
两份 $2x\cos x$ 要合并成 $4x\cos x$。对 $z=x\ln x$（$x>0$），$z'=\ln x+1$，而不是 $1/x$。
\end{example}

\begin{example}{三个因子：逐项求导而不遗漏}{基础训练}
令 $y=xe^x\sin x$，则
\begin{align*}
 y'&=e^x\sin x+xe^x\sin x+xe^x\cos x\\
   &=e^x\bigl[(x+1)\sin x+x\cos x\bigr].
\end{align*}
第一行对应三个因子分别求导。也可把 $e^x\sin x$ 作为整体，先用两因子法则，再在整体内部用一次乘积法则；两种分组结果相同。
\end{example}

\paragraph{乘积内部还可能有复合关系}
若 $y=u(x)v(x)$，乘积法则中的 $u',v'$ 仍须各自正确计算。例如
\[
 \bigl[e^{x^2}\sin(3x)\bigr]'
 =2xe^{x^2}\sin(3x)+3e^{x^2}\cos(3x).
\]
先分清“两个因子的乘积”，再处理各因子内部的复合关系，不能漏掉内层的 $2x$ 或 $3$。若是 $[u(x)]^m$，可用 $m u^{m-1}u'$；它不是 $m$ 个不同函数的导数相乘。

\subsection{高阶乘法求导：二项式系数从哪里来}\label{r2:product:higher}
\begin{align*}
 (uv)''&=u''v+2u'v'+uv'',\\
 (uv)'''&=u'''v+3u''v'+3u'v''+uv''',\\
 (uv)^{(n)}&=\sum_{k=0}^{n}\binom nk u^{(k)}v^{(n-k)}\quad(n\ge0).
\end{align*}
这里约定 $u^{(0)}=u$，$\binom nk=n!/[k!(n-k)!]$。每个加项的总求导阶数都是 $n$；两端项系数都是一。二阶的中间系数二、三阶的中间系数三，都是再次求导时相同项合并的结果。最后一式即 Leibniz 公式，证明见第\ref{sec:2:2}节。

三个或更多因子的高阶推广为
\[
 \left(\prod_{j=1}^m u_j\right)^{(n)}
 =\sum_{k_1+\cdots+k_m=n}
   \frac{n!}{k_1!\cdots k_m!}\prod_{j=1}^m u_j^{(k_j)},
 \qquad k_j\ge0.
\]
求和遍历所有非负整数阶数组合。例如三因子二阶展开有三项“一个因子求二阶”和三项系数为二的“两个因子各求一阶”：
\begin{align*}
 (uvw)''={}&u''vw+uv''w+uvw''\\
 &+2u'v'w+2u'vw'+2uv'w'.
\end{align*}

\begin{method}{高阶乘积求导的省算顺序}
若其中一个因子是低次多项式，优先让它在 Leibniz 和式中承受 $k$ 阶导数，超过次数的项自动为零。另一个因子若为指数或三角函数，先写其通项。最后用 $n=1,2$ 的直接求导核对系数，不用“把两边各求 $n$ 阶再相乘”的错误规则。
\end{method}

\begin{example}{多项式乘一般函数}{综合训练}
当 $f$ 有所需导数时，
\[
 (xf)^{(n)}=xf^{(n)}+nf^{(n-1)}\quad(n\ge1),
\]
\[
 (x^2f)^{(n)}=x^2f^{(n)}+2nxf^{(n-1)}+n(n-1)f^{(n-2)}\quad(n\ge2).
\]
第二式来自 $k=0,1,2$ 三项，因 $\binom n2(x^2)''=n(n-1)$。例如 $f=e^{ax}$ 时，各阶导数可保留为 $(e^{ax})^{(j)}$ 再代入，避免在 $a=0$ 时机械写出负幂。前文的 $x^2e^x$ 是 $a=1$ 的特例。
\end{example}

\paragraph{多因子对数写法：只在非零区间使用}
若所有因子非零，$y=\prod u_j$ 可写
\[
 \frac{y'}y=\sum_j\frac{u_j'}{u_j},\qquad
 y'=y\sum_j\frac{u_j'}{u_j}.
\]
这是对 $\ln|y|=\sum_j\ln|u_j|$ 求导，不允许在任一因子为零的点直接除法。零点处应回到原来的逐项求导式；若通过化简得到连续延拓，也须说明其来源。

\paragraph{自测与完整答案}
\begin{enumerate}
\item 求 $x^2\cos x$ 的导数。答：$2x\cos x-x^2\sin x$。
\item 求 $(1+x)e^{2x}$ 的二阶导数。答：$4(x+2)e^{2x}$；一阶为 $(2x+3)e^{2x}$，再求导可核对。
\item 求 $x\sin x\cos x$ 的导数。答：$\sin x\cos x+x(\cos^2x-\sin^2x)$；也可先化为 $\tfrac{x}{2}\sin2x$。
\item 写出 $(x^3f)^{(n)}$，$n\ge3$。答：$x^3f^{(n)}+3nx^2f^{(n-1)}+3n(n-1)xf^{(n-2)}+n(n-1)(n-2)f^{(n-3)}$。
\end{enumerate}

\subsection{复合函数求导公式全表：外导数乘内导数}\label{r2:chain}
本节与第\ref{r2:product}节的乘法求导配合使用。“乘积”指多个因子相乘；“复合”指一个函数作为另一个函数的输入，同一个表达式可能同时含有这两种结构。

\begin{theorem}{两层、多层复合的一阶及二三阶公式}
若 $u$ 在所求点可导，$f$ 在 $u(x)$ 处可导，则
\[
 [f(u(x))]'=f'(u(x))u'(x).
\]
若 $y=f(g(h(x)))$，则
\[
 y'=f'(g(h(x)))g'(h(x))h'(x).
\]
各层均需在对应取值处可导。当所需高阶导数存在时，
\begin{align*}
 [f(u)]''&=f''(u)(u')^2+f'(u)u'',\\
 [f(u)]'''&=f'''(u)(u')^3+3f''(u)u'u''+f'(u)u'''.
\end{align*}
若内层为 $ax+b$，则 $[f(ax+b)]^{(n)}=a^n f^{(n)}(ax+b)$（$n\ge1$）。一般非线性内层不能照搬这个简化式。
\end{theorem}
\begin{proof}
设 $\Delta u=u(x+h)-u(x)$，外层可导给 $\Delta f=f'(u)\Delta u+\varepsilon(\Delta u)\Delta u$，其中 $\varepsilon(s)\to0$，并可设 $\varepsilon(0)=0$。内层可导给 $\Delta u/h\to u'$，故除以 $h$ 后余项趋零，即得链式法则。这种证明不需要除以可能为零的 $\Delta u$。二、三阶式对上一阶表达式整体使用乘积与链式法则即可推出。
\end{proof}

以下统一记 $u=u(x),v=v(x)$，$a,\alpha$ 为常数，表中公式只在实际定义域内且满足所列条件时使用。令 $u=x,u'=1$ 就恢复基本初等函数导数表。
\begin{longtable}{@{}p{.23\textwidth}p{.40\textwidth}p{.31\textwidth}@{}}
\toprule 函数&关于 $x$ 的导数&实值区间与条件\\\midrule\endhead
$[u(x)]^\alpha$&$\alpha u^{\alpha-1}u'$&任意实 $\alpha$ 统一取 $u>0$\\
$1/u$&$-u'/u^2$&$u\ne0$\\
$\sqrt{u}$&$u'/(2\sqrt{u})$&此公式要求 $u>0$\\
$e^u$&$e^u u'$&$u$ 可导\\
$a^u$&$a^u\ln a\,u'$&$a>0$；$a=1$ 时导数零\\
$\ln u$&$u'/u$&$u>0$\\
$\ln|u|$&$u'/u$&$u\ne0$\\
$\log_a u$&$u'/(u\ln a)$&$u>0,a>0,a\ne1$\\
$\sin u$&$\cos u\,u'$&弧度制\\
$\cos u$&$-\sin u\,u'$&注意负号\\
$\tan u$&$\sec^2u\,u'$&$\cos u\ne0$\\
$\cot u$&$-\csc^2u\,u'$&$\sin u\ne0$\\
$\sec u$&$\sec u\tan u\,u'$&$\cos u\ne0$\\
$\csc u$&$-\csc u\cot u\,u'$&$\sin u\ne0$\\
$\arcsin u$&$u'/\sqrt{1-u^2}$&$|u|<1$\\
$\arccos u$&$-u'/\sqrt{1-u^2}$&$|u|<1$\\
$\arctan u$&$u'/(1+u^2)$&$u$ 可导\\
$\operatorname{arccot}u$&$-u'/(1+u^2)$&取值为 $(0,\pi)$ 的分支\\
$u^v$&$u^v\bigl(v'\ln u+vu'/u\bigr)$&统一取 $u>0$，$u,v$ 可导\\
\bottomrule
\end{longtable}

对于正整数幂，幂函数公式可以在全部实值点使用；负整数幂仍排除零。表中根式或反三角函数的内部条件失效时，应回到复合函数本身或导数定义判断：$\sqrt{x^2}=|x|$ 在零不可导，$\sqrt{x^4}=x^2$ 却可导，不能仅因表格出现零分母就否定所有特殊复合的可导性。

\paragraph{可选的反三角与双曲补表}
若采用 $\operatorname{arcsec}u=\arccos(1/u)$、$\operatorname{arccsc}u=\arcsin(1/u)$ 的分支约定，则在 $|u|>1$ 上
\[
 (\operatorname{arcsec}u)'=\frac{u'}{|u|\sqrt{u^2-1}},\qquad
 (\operatorname{arccsc}u)'=-\frac{u'}{|u|\sqrt{u^2-1}}.
\]
这两个函数的命名与分支并非所有教材统一，使用前应先核对约定。双曲函数的常见复合形式为
\begin{align*}
 (\sinh u)'&=\cosh u\,u',&(\cosh u)'&=\sinh u\,u',\\
 (\tanh u)'&=\frac{u'}{\cosh^2u},&(\coth u)'&=-\frac{u'}{\sinh^2u}\quad(u\ne0).
\end{align*}
这张可选补表不改变各卷种的考试范围。

\begin{example}{三层复合与二阶导数}{基础训练}
设 $y=\ln(1+\sin^2(3x))$。由外向内依次求导，
\[
 y'=\frac{1}{1+\sin^2(3x)}\,2\sin(3x)\cos(3x)\,3
 =\frac{6\sin(3x)\cos(3x)}{1+\sin^2(3x)}.
\]
分母至少为一，全实轴合法。另取 $z=\sin(x^2)$，则 $z'=2x\cos(x^2)$，二阶为 $z''=2\cos(x^2)-4x^2\sin(x^2)$；只保留外层二阶项会漏掉 $2\cos(x^2)$。
\end{example}

\begin{example}{变底变指数：乘积与复合一起使用}{综合训练}
对 $y=(1+x^2)^{\sin x}$，底数恒正，取 $\ln y=\sin x\ln(1+x^2)$。右边先用乘积法则，左边用链式法则，得到
\[
 y'=(1+x^2)^{\sin x}\left[\cos x\ln(1+x^2)+\frac{2x\sin x}{1+x^2}\right].
\]
对 $y=x^x$（$x>0$），同理得 $y'=x^x(\ln x+1)$。当指数本身变化时，不能只套常数指数幂法则。
\end{example}

\paragraph{隐式、参数与反函数的相邻公式}
隐函数 $F(x,y)=0$ 在 $F\in C^2$、基点满足方程且 $F_y\ne0$ 时，
\[
 y'=-\frac{F_x}{F_y},\qquad
 y''=-\frac{F_{xx}+2F_{xy}y'+F_{yy}(y')^2}{F_y}.
\]
反函数 $g=f^{-1}$ 的一阶公式为 $g'(f(a))=1/f'(a)$，需局部可逆且 $f'(a)\ne0$；若二阶条件满足，则 $g''(f(a))=-f''(a)/[f'(a)]^3$。反函数不是倒数，不能把 $f^{-1}(x)$ 与 $1/f(x)$ 混为一谈。参数求导仍需用 $\dd/\dd x=[1/x'(t)]\dd/\dd t$，先核对 $x'(t)\ne0$。

\paragraph{自测与答案}
\begin{enumerate}
\item $\arctan(e^{2x})$：导数为 $2e^{2x}/(1+e^{4x})$。
\item $\sqrt{\ln x}$：在 $x>1$ 的导数为 $1/(2x\sqrt{\ln x})$，不把内部公式直接代到端点 $x=1$。
\item $\log_2(1+x^2)$：导数为 $2x/[(1+x^2)\ln2]$。
\item $x^{\sin x}$：在 $x>0$ 的导数为 $x^{\sin x}[\cos x\ln x+\sin x/x]$。
\end{enumerate}
链式结构可对照 \href{https://openstax.org/books/calculus-volume-1/pages/3-6-the-chain-rule}{OpenStax 微积分教材的链式法则章节}；本表将内函数及适用条件显式写出。


\subsection{例题与方法比较}
\begin{example}{基础例题｜对数求导}{难度：$\star$}
\label{ex:2:1}\textbf{题目} 求$x^x$的导数，$x>0$。\par
\textbf{思路} 底数和指数都变化。\par
\textbf{解答} 方法一：设$y=x^x$，$\ln y=x\ln x$，得$y\prime=x^x(\ln x+1)$。方法二：写成$\e^{x\ln x}$再用链式法则，得相同结果。\par
\textbf{方法点评} 对数求导便于多因子，指数表示便于高阶推广。\par
\textbf{易错点或变式} 不能用常数指数幂法则直接得到$xx^{x-1}$。
\end{example}
\begin{example}{考研风格例题｜隐函数}{难度：$\star\star$}
\label{ex:2:2}\textbf{题目} 在单位圆上$x y\ne0$处求$y\prime,y\prime\prime$。\par
\textbf{思路} 令 $F=x^2+y^2-1$，基点在圆上且 $F\in C^\infty$，题设 $y\ne0$ 给 $F_y=2y\ne0$，故存在局部光滑分支。先求一阶，再对含 $y$ 的表达式求导。\par
\textbf{解答} 方法一：$x^2+y^2=1$给$y\prime=-\frac{x}{y}$，再求导得$y\prime\prime=-\frac{1}{y}-\frac{x^2}{y^3}=-\frac{1}{y^3}$。方法二：局部分支$y=\pm\sqrt{1-x^2}$，直接求两次导数得到同式。\par
\textbf{方法点评} 隐式法统一上下分支，显式法便于核对。\par
\textbf{易错点或变式} 在$y=0$处不能把圆写成局部可导的$y(x)$。
\end{example}
\begin{example}{考研风格例题｜参数曲线}{难度：$\star\star$}
\label{ex:2:3}\textbf{题目} $x=t^2,y=t^3$，$t\ne0$，求二阶导数。\par
\textbf{思路} 第二次仍要除以$x\prime(t)$。\par
\textbf{解答} 方法一：$\frac{\dd y}{\dd x}=\frac{3t}{2}$，$\frac{\dd^2y}{\dd x^2}=\frac{\frac{3}{2}}{2t}=\frac{3}{4t}$。方法二：在$t>0$与$t<0$分别写$y=\pm x^{\frac{3}{2}}$，求导后代$x=t^2$。\par
\textbf{方法点评} 参数法无需分支讨论，消参法检查符号。\par
\textbf{易错点或变式} $t=0$是奇异参数点，公式不能直接套用。
\end{example}
\section{高阶导数与Leibniz公式}\label{sec:2:2}

\begin{theorem}{常用导数公式，标准结论}
在各自可导的实值区间上，
\begin{align*}
(x^\alpha)\prime&=\alpha x^{\alpha-1},& (\ln|x|)\prime&=\dfrac{1}{x},\\
(a^x)\prime&=a^x\ln a\quad(a>0),& (\e^x)\prime&=\e^x,\\
(\tan x)\prime&=\sec^2x,&(\cot x)\prime&=-\csc^2x,\\
(\arcsin x)\prime&=(1-x^2)^{-\frac{1}{2}},&(\arccos x)\prime&=-(1-x^2)^{-\frac{1}{2}},\\
(\arctan x)\prime&=(1+x^2)^{-1}.&&
\end{align*}
\end{theorem}
\begin{proof}幂、指数和对数公式可由重要极限与导数定义推出；反三角公式由反函数求导获得。例如$y=\arcsin x$时$\sin y=x$且内部$\cos y>0$，所以$y\prime=\frac{1}{\cos y}=\frac{1}{\sqrt{1-x^2}}$。端点$\pm1$不能套内部导数公式。\end{proof}

\subsection{概念、定理与推导}

\begin{theorem}{Leibniz公式}若$u,v$有$n$阶导数，则
\begin{equation}(uv)^{(n)}=\sum_{k=0}^{n}\binom nk u^{(k)}v^{(n-k)}.\end{equation}
\end{theorem}
\begin{proof}对$n$归纳。再求导时相邻两项系数相加，由$\binom nk+\binom n{k-1}=\binom{n+1}k$完成归纳。\end{proof}
常用结构：$(\e^{ax})^{(n)}=a^n\e^{ax}$，$(\sin ax)^{(n)}=a^n\sin(ax+\frac{n\pi}{2})$，$(\frac{1}{x})^{(n)}=\frac{(-1)^n n!}{x^{n+1}}$。高阶微分一般不具备一阶微分形式不变性，因为$d^2y=f\prime\prime(x)(dx)^2+f\prime(x)d^2x$。

\begin{method}{高阶导数与Leibniz公式}条件：乘积两因子足够光滑。步骤：优先让多项式因子参与高阶导数以截断求和；用部分分式分解有理函数。易错点：阶乘和二项式系数漏写。\end{method}
\subsection{例题与方法比较}
\begin{example}{考研风格例题｜多项式乘指数}{难度：$\star\star$}
\label{ex:2:4}\textbf{题目} 求$(x^2\e^x)^{(n)}$，$n\ge2$。\par
\textbf{思路} 多项式三阶导数起为零。\par
\textbf{解答} 方法一：Leibniz给$\e^x[x^2+2nx+n(n-1)]$。方法二：使用$(D+1)^n x^2$，仅零、一、二阶$D$项非零，得到同式。\par
\textbf{方法点评} Leibniz透明，移位算子便于重复应用。\par
\textbf{易错点或变式} 二阶项是$\binom n2\cdot2=n(n-1)$。
\end{example}
\begin{example}{考研风格例题｜分式高阶}{难度：$\star\star$}
\label{ex:2:5}\textbf{题目} 求$f(x)=\frac{1}{x^2-1}$的$n$阶导数。\par
\textbf{思路} 把分式拆成一次因子的倒数。\par
\textbf{解答} 方法一：$f=\frac12[(x-1)^{-1}-(x+1)^{-1}]$，故$f^{(n)}=\frac{(-1)^n n!}{2}[(x-1)^{-n-1}-(x+1)^{-n-1}]$。方法二：先求$n=0$式再归纳求导，系数递乘$-(n+1)$，得同式。\par
\textbf{方法点评} 部分分式找结构，归纳证明对所有阶成立。\par
\textbf{易错点或变式} 定义域排除$\pm1$。
\end{example}
\begin{example}{基础例题｜振荡导数}{难度：$\star$}
\label{ex:2:6}\textbf{题目} 求$\sin2x$的第十阶导数。\par
\textbf{思路} 三角函数导数以四为周期。\par
\textbf{解答} $2^{10}\sin(2x+5\pi)=-1024\sin2x$。可先算二阶为$-4\sin2x$，再重复五次。\par
\textbf{方法点评} 周期公式减少机械计算。\par
\textbf{易错点或变式} 内层系数2要出现十次。
\end{example}
\section{中值定理与辅助函数}\label{sec:2:3}

\begin{method}{辅助函数构造的三个入口}
目标含$f\prime-\lambda f$，逆用乘积法则构造$\e^{-\lambda x}f$；目标含端点割线斜率，减去通过两端点的直线；目标涉及二阶导数和两个端点，减去插值直线及待定倍数$(x-a)(x-b)$。这三种入口分别对应一阶线性表达式、弦与切线比较、二阶误差控制。构造后必须核验连续、可导和零点条件。
\end{method}

\subsection{概念、定理与推导}

\begin{theorem}{Fermat、Rolle与Lagrange}内点可导极值处$f\prime=0$。若$f\in C[a,b]$、在$(a,b)$可导且$f(a)=f(b)$，存在$\xi$使$f\prime(\xi)=0$。去掉端点相等条件，则存在$\xi$使$f(b)-f(a)=f\prime(\xi)(b-a)$。\end{theorem}
\begin{proof}Fermat由左右差商异号得出。Rolle利用连续函数取最值：若非恒定，至少一项极值在内点。Lagrange对$f(x)-f(a)-\frac{[f(b)-f(a)](x-a)}{b-a}$应用Rolle。\end{proof}
\begin{theorem}{Cauchy中值定理}若$f,g$连续于$[a,b]$且内部可导，存在$\xi$使$[f(b)-f(a)]g\prime(\xi)=[g(b)-g(a)]f\prime(\xi)$。若要写成比式，须另外保证所除的量非零。\end{theorem}
\begin{proof}对$[f(b)-f(a)]g(x)-[g(b)-g(a)]f(x)$使用Rolle。\end{proof}

\begin{method}{中值定理与辅助函数}条件：连续闭区间、可导开区间。步骤：目标是$f\prime-\lambda f$时乘$\e^{-\lambda x}$；目标含线性弦斜率时减弦。易错点：不同次中值定理所得中值点通常不同。\end{method}
\subsection{例题与方法比较}
\begin{example}{考研风格例题｜指数辅助}{难度：$\star\star$}
\label{ex:2:7}\textbf{题目} $f(0)=f(1)=0$，$f$连续且内部可导，证明某$\xi$满足$f\prime(\xi)=f(\xi)$。\par
\textbf{思路} 寻找导数中出现$f\prime-f$的乘积。\par
\textbf{解答} 令$F(x)=\e^{-x}f(x)$，两端为零，Rolle给$F\prime(\xi)=\e^{-\xi}[f\prime(\xi)-f(\xi)]=0$。另法：把目标看成一阶线性表达式，用积分因子$\e^{-x}$重建同一辅助函数。\par
\textbf{方法点评} 积分因子思路解释辅助函数的构造；核心定理仍为Rolle。\par
\textbf{易错点或变式} 不要求$f\prime$连续。
\end{example}
\begin{example}{考研风格例题｜对数不等式}{难度：$\star\star$}
\label{ex:2:8}\textbf{题目} 证明$\frac{x}{1+x}<\ln(1+x)<x$，$x>0$。\par
\textbf{思路} 导数、积分或中值定理均可。\par
\textbf{解答} 方法一：$\ln(1+x)=\frac{x}{1+\xi}$，$0<\xi<x$，直接夹逼。方法二：在$[0,x]$上积分$\frac{1}{1+x}\le\frac{1}{1+t}\le1$，由于非恒定得到严格不等式。方法三：两差函数导数分别为$\frac{x}{(1+x)^2}$与$\frac{x}{1+x}$，从0递增。\par
\textbf{方法点评} 中值法简洁，积分法显示面积，导数法适合复杂参数。\par
\textbf{易错点或变式} $x<0$时不可照搬未经分析的不等号方向。
\end{example}
\begin{example}{考研风格例题｜零点个数}{难度：$\star\star$}
\label{ex:2:9}\textbf{题目} 证明$x^3+x-1=0$只有一个实根。\par
\textbf{思路} 先存在再唯一。\par
\textbf{解答} $f(0)=-1,f(1)=1$，连续性保证根。$f\prime=3x^2+1>0$，严格递增保证唯一；也可假设有两根，用Rolle推出导数零点而矛盾。\par
\textbf{方法点评} 单调法与Rolle反证都是全局唯一性工具。\par
\textbf{易错点或变式} 只有导数为正不能单独保证根存在。
\end{example}
\section{Taylor、极值、凹凸与曲率}\label{sec:2:4}
\subsection{概念、定理与推导}

\begin{definition}{极值与凸性}局部极小值要求定义域内$a$的某邻域中$f(x)\ge f(a)$；若对$x\ne a$均严格大于，称严格局部极小值。设$I$为区间，凸函数满足对任意$x,y\in I$和$0\le t\le1$有$f((1-t)x+ty)\le(1-t)f(x)+tf(y)$；在开区间上二阶可导时，凸等价于$f\prime\prime\ge0$处处成立。凹函数取相反的不等号。函数图像在某连续点两侧的凹凸性改变时，该点称为拐点；涉及竖切线时须说明所采用的约定。\end{definition}
\begin{warning}本册采用凸分析的命名：凸函数的图像位于弦下，二阶可导时$f\prime\prime\ge0$。部分高等数学教材把这种图形称为“凹弧”或“凹向上”，并把$f\prime\prime\le0$的弧称为“凸弧”。函数的凸性与弧的凹凸称呼可能相反，判断时应以弦的不等式和二阶导数符号为准。\end{warning}
\begin{theorem}{判别与几何公式}驻点$f\prime(a)=0$处，若$f\prime\prime(a)>0$则严格极小，若小于0则严格极大；等于0须进一步判断。图像曲率
\begin{equation}\kappa=\dfrac{|f\prime\prime|}{(1+(f\prime)^2)^{\frac{3}{2}}},\qquad \rho=\dfrac{1}{\kappa}\quad(\kappa>0).\end{equation}
\end{theorem}
\begin{proof}二阶判别来自$f(a+h)-f(a)=\frac{f\prime\prime(a)h^2}{2}+o(h^2)$。曲率定义为切向角对弧长的变化率，利用$\tan\theta=f\prime$和$\frac{\dd s}{\dd x}=\sqrt{1+(f\prime)^2}$得到公式。\end{proof}
斜渐近线$y=kx+b$需$k=\lim \frac{f(x)}{x}$及$b=\lim(f(x)-kx)$均存在；竖渐近线由单侧无穷极限判断。

\begin{method}{Taylor、极值、凹凸与曲率}条件：区间内分段光滑。步骤：列驻点、不可导点与端点；用导数符号表判单调和凹凸。易错点：$f\prime\prime=0$只是候选拐点。\end{method}
\subsection{例题与方法比较}
\begin{example}{基础例题｜极值}{难度：$\star$}
\label{ex:2:10}\textbf{题目} 分析$f=x^3-3x$的单调性与极值。\par
\textbf{思路} 一阶导数符号分区间。\par
\textbf{解答} $f\prime=3(x^2-1)$，在$(-\infty,-1)$与$(1,\infty)$递增，在$(-1,1)$递减。$f(-1)=2$为极大，$f(1)=-2$为极小。$f\prime\prime=6x$在0变号，原点为拐点。\par
\textbf{方法点评} 符号表同时呈现局部与全局形状。\par
\textbf{易错点或变式} 在全实数轴没有最大值和最小值。
\end{example}
\begin{example}{考研风格例题｜曲率}{难度：$\star\star$}
\label{ex:2:11}\textbf{题目} 求抛物线$y=x^2$在原点的曲率半径。\par
\textbf{思路} 代入曲率公式。\par
\textbf{解答} $f\prime=2x,f\prime\prime=2$，$\kappa(0)=2$，故$\rho=\frac{1}{2}$。密切圆圆心在法线正向$(0,\frac{1}{2})$。\par
\textbf{方法点评} 曲率反映切线方向的局部变化速度。\par
\textbf{易错点或变式} 曲率半径与函数值或切线斜率不是同一量。
\end{example}
\begin{example}{考研风格例题｜退化驻点}{难度：$\star\star$}
\label{ex:2:12}\textbf{题目} 比较$x^4$与$x^3$在0的极值和拐点。\par
\textbf{思路} 二阶导数在0均为零，必须看符号。\par
\textbf{解答} $x^4\ge0$，0是严格极小值；其二阶导数$12x^2$不变号，无拐点。$x^3$左右异号，无极值；二阶导数$6x$变号，有拐点。\par
\textbf{方法点评} 直接比较函数值和二阶符号比高阶公式更稳。\par
\textbf{易错点或变式} 二阶导数为零不能判定任何一种结论。
\end{example}
\section{中值构造与不等式拓展}\label{sec:2:5}
\subsection{概念、定理与推导}

\begin{lemma}{有界导数的Lipschitz估计}若区间上$|f\prime|\le M$，则$|f(x)-f(y)|\le M|x-y|$。\end{lemma}
\begin{proof}对两点之间使用Lagrange中值定理再取绝对值。\end{proof}
\begin{extension}二阶不等式常用“减去二次多项式”，一阶线性导数组合常用积分因子，多个零点常用反复Rolle。寻找辅助函数时可先对目标表达式逆向积分。\end{extension}

\begin{method}{中值构造与不等式拓展}条件：需要的导数存在且端点条件足以应用中值定理。步骤：把结论整理为某函数导数为零或非负，再构造该函数。易错点：只写辅助函数不验证端点与可导性。\end{method}
\subsection{例题与方法比较}
\begin{example}{竞赛风格例题｜二阶估计}{难度：$\star\star\star$}
\label{ex:2:13}\textbf{题目} $f\in C^2[0,1]$，$f(0)=f(1)=0$且$|f\prime\prime|\le M$，证$|f(x)|\le \frac{Mx(1-x)}{2}$。\par
\textbf{思路} 用二次函数包络。\par
\textbf{解答} 设$u=f-\frac{Mx(1-x)}{2}$，则$u\prime\prime=f\prime\prime+M\ge0$且两端为零；凸函数在弦下，故$u\le0$。同理$v=f+\frac{Mx(1-x)}{2}$凹且两端零，$v\ge0$。合并得结论。\par
\textbf{方法点评} 比较曲率可产生全区间估计。\par
\textbf{易错点或变式} 上包络对应的差函数是凸函数，符号不能反。
\end{example}
\begin{example}{竞赛风格例题｜插值余项}{难度：$\star\star\star$}
\label{ex:2:14}\textbf{题目} $f\in C^2[a,b]$，$\ell$为端点线性插值，证明内点$x$处$f(x)-\ell(x)=\frac{f\prime\prime(\xi)(x-a)(x-b)}{2}$。\par
\textbf{思路} 添加二次因子制造三个零点。\par
\textbf{解答} 令$K=\frac{f(x)-\ell(x)}{(x-a)(x-b)}$，$F(t)=f(t)-\ell(t)-K(t-a)(t-b)$在$a,x,b$为零。两次Rolle得某$\xi$使$F\prime\prime(\xi)=f\prime\prime(\xi)-2K=0$，代回即得。\par
\textbf{方法点评} 待定系数法把误差表示变成零点问题。\par
\textbf{易错点或变式} $\xi$依赖于所选$x$。
\end{example}
\begin{example}{竞赛风格例题｜指数切线}{难度：$\star\star\star$}
\label{ex:2:15}\textbf{题目} 证明$\e^x\ge1+x$且等号仅在0。\par
\textbf{思路} 凸性或最小值。\par
\textbf{解答} 令$g=\e^x-1-x$，$g\prime=\e^x-1$，在负半轴负、正半轴正，唯一最小值$g(0)=0$。另一法：指数严格凸，图像在0处切线$1+x$上方。\par
\textbf{方法点评} 导数法自足，凸性法适合推广到Jensen不等式。\par
\textbf{易错点或变式} 不要将局部Taylor近似未经余项控制当成全局证明。
\end{example}
\begin{warning}可导推出连续，反向不成立；驻点不一定是极值点；拐点必须检查凹凸是否变号；高阶微分不能直接照搬一阶形式不变性。\end{warning}


\section{微分中的存在性、插值与误差控制}\label{mat:differentiation}
\begin{examtip}补足Darboux性质、高阶中值构造、导数估计、渐近线与近似计算。\end{examtip}
\subsection{导数的Darboux性质}\label{mat:differentiation:darboux}
\begin{method}{导数的Darboux性质}\textbf{适用条件} 函数在区间上可导，导数不必连续。\par\textbf{实施与依据} 对$f(x)-\lambda x$取内点极值，证明导数具有介值性。\end{method}
\begin{example}{专题例题｜导数的Darboux性质}{难度：$\star\star$}
\textbf{题目} 设$f$在包含$[a,b]$的开区间上可导，$a<b$。若$f\prime(a)<\lambda<f\prime(b)$，证明某$c\in(a,b)$满足$f\prime(c)=\lambda$。\par\textbf{思路与解答} 令$g=f-\lambda x$，有$g\prime(a)<0,g\prime(b)>0$，故$g$在两端都不能取得最小值。连续性保证在闭区间取最小，最小点必在内部，Fermat给$g\prime(c)=0$。\par\textbf{方法点评与边界} 导数可振荡不连续，但不能发生跳跃间断。
\end{example}
\subsection{反复Rolle与零点重数}\label{mat:differentiation:rolle}
\begin{method}{反复Rolle与零点重数}\textbf{适用条件} 多个零点或多阶导数关系。\par\textbf{实施与依据} 相邻零点间产生导数零点，反复次数与所需阶数对应。\end{method}
\begin{example}{专题例题｜反复Rolle与零点重数}{难度：$\star\star$}
\textbf{题目} 证明非零$n$次多项式至多有$n$个互异实根。\par\textbf{思路与解答} 若有$n+1$个根，反复Rolle推出$n$阶导数至少有一零点；但$n$阶导数是非零常数，矛盾。\par\textbf{方法点评与边界} 重根问题要将重数与导数消失一并计算。
\end{example}
\subsection{Lagrange插值误差}\label{mat:differentiation:interpolate}
\begin{method}{Lagrange插值误差}\textbf{适用条件} $n+1$个互异节点，函数有$n+1$阶连续导数。\par\textbf{实施与依据} 减去插值多项式及$K\prod(x-x_j)$，制造$n+2$个零点。\end{method}
\begin{example}{专题例题｜Lagrange插值误差}{难度：$\star\star$}
\textbf{题目} 写出插值余项。\par\textbf{思路与解答} 设$p_n(x_j)=f(x_j)$，对非节点$x$令$K=\frac{f(x)-p_n(x)}{\prod_j(x-x_j)}$。反复Rolle得到$K=\frac{f^{(n+1)}(\xi)}{(n+1)!}$，故误差等于该系数乘节点乘积。\par\textbf{方法点评与边界} 节点误差为0；估计最大误差还需控制导数和节点乘积。
\end{example}
\subsection{渐近线与完整函数分析}\label{mat:differentiation:asymptote}
\begin{method}{渐近线与完整函数分析}\textbf{适用条件} 端点、无穷远与奇点。\par\textbf{实施与依据} 先定义域，再查竖、水平和斜渐近线，再列单调与凹凸表。\end{method}
\begin{example}{专题例题｜渐近线与完整函数分析}{难度：$\star\star$}
\textbf{题目} 求$f(x)=x+\frac{1}{x}$的渐近线。\par\textbf{思路与解答} $x\to0^\pm$时$\frac{1}{x}$无穷，故$x=0$为竖渐近线。$x\to\pm\infty$时$f(x)-x\to0$，斜渐近线$y=x$。\par\textbf{方法点评与边界} 渐近线允许曲线在有限位置穿过，并不要求始终位于同侧。
\end{example}
\subsection{Newton法与局部收敛}\label{mat:differentiation:newton}
\begin{method}{Newton法与局部收敛}\textbf{适用条件} $f$在根$r$附近属于$C^2$，$f(r)=0,f\prime(r)\ne0$，且初值充分接近$r$。\par\textbf{实施与依据} 由连续性缩小邻域，使$|f\prime|$有正下界且$|f\prime\prime|$有上界。迭代$x_{n+1}=x_n-\frac{f(x_n)}{f\prime(x_n)}$；Taylor给$|x_{n+1}-r|\le C|x_n-r|^2$。将初始误差取到$C|x_1-r|<1$且迭代留在该邻域内，即得局部二次收敛。\end{method}
\begin{example}{专题例题｜Newton法与局部收敛}{难度：$\star\star$}
\textbf{题目} 求$\sqrt a$的Newton迭代，$a>0$。\par\textbf{思路与解答} 对$x^2-a$迭代得$x_{n+1}=\frac{x_n+\frac{a}{x_n}}{2}$。若$x_n>0$，下一项不小于$\sqrt a$；并有$x_{n+1}-\sqrt a=\frac{(x_n-\sqrt a)^2}{2x_n}$，说明接近根后二次收敛。\par\textbf{方法点评与边界} 导数接近0或初值远离根时不能保证全局收敛。
\end{example}




\section{微分、中值定理与不等式的构造训练}\label{deep:differential_source_deep}

\begin{examtip}补齐高阶导数、隐参数求导、辅助函数构造、Taylor余项、曲率和渐近线之间的联系。\end{examtip}

\subsection{复合、隐函数、参数与高阶导数}\label{deep:differential_source_deep:derivatives}

\begin{method}{复合、隐函数、参数与高阶导数}\textbf{适用条件} 相关函数在所考察点可微，隐式求导的分母非零。\par\textbf{步骤与依据} 先写依赖关系；二阶求导把一阶结果整体再求导；参数方程使用$\frac{\dd}{\dd x}=\frac{1}{x'(t)}\frac{\dd}{\dd t}$。\par\textbf{易错点} 二阶导数不是$\frac{y''(t)}{x''(t)}$，隐函数中$y$随$x$变化。\end{method}

\begin{example}{考研风格｜对数求导与Taylor互证}{难度：$\star\star\star$}\label{deep:ex:differential_source_deep:logdiff}\textbf{题目} 求$y=x^x$（$x>0$）的一、二阶导数，并给$x=1$处二阶Taylor式。\par\textbf{思路} 先取对数，或写成指数复合。\par\textbf{详细解答} $\ln y=x\ln x$，故$y'=x^x(1+\ln x)$，$y''=x^x[(1+\ln x)^2+\frac{1}{x}]$。在1处$y=1,y'=1,y''=2$，所以$x^x=1+(x-1)+(x-1)^2+o((x-1)^2)$。另一法令$h=x-1$，$(1+h)\ln(1+h)=h+\frac{h^2}{2}+o(h^2)$，再指数展开得到同式。\par\textbf{方法点评} 对数求导适合变底变指数；局部展开可核对导数。\par\textbf{变式与易错点} 不能写$(x^x)'=x\,x^{x-1}$，那只适用于常指数。\end{example}

\begin{example}{考研风格｜隐函数二阶导数}{难度：$\star\star\star$}\label{deep:ex:differential_source_deep:implicit2}\textbf{题目} 设$x^2+xy+y^2=3$，求过$(1,1)$的局部支在该点的$y',y''$。\par\textbf{思路} 先确认$F_y=x+2y$非零。\par\textbf{详细解答} $F_y(1,1)=3$，存在局部可微函数。一次微分$2x+y+(x+2y)y'=0$，得$y'(1)=-1$。再微分$2+2y'+2(y')^2+(x+2y)y''=0$，代入得$y''(1)=-\frac{2}{3}$。另一法在点附近令$y=1+ah+bh^2+o(h^2),x=1+h$，比较系数得$a=-1,b=-\frac{1}{3}$，所以$y''=2b=-\frac{2}{3}$。\par\textbf{方法点评} 局部展开法适合只求某点导数，避免解完整二次根式。\par\textbf{变式与易错点} 把$xy'$再求导时会产生$y'$，不能漏乘积项。\end{example}

\begin{example}{考研风格｜摆线曲率与参数导数}{难度：$\star\star\star$}\label{deep:ex:differential_source_deep:parametric}\textbf{题目} 摆线$x=a(t-\sin t),y=a(1-\cos t)$，$a>0,0<t<2\pi$。求$\frac{\dd y}{\dd x}$和曲率。\par\textbf{思路} 参数导数公式与平面曲率行列式。\par\textbf{详细解答} $x'=a(1-\cos t),y'=a\sin t$，故$y_x'=\frac{\sin t}{1-\cos t}=\cot(\frac{t}{2})$。曲率$\kappa=\frac{|x'y''-y'x''|}{(x'^2+y'^2)^{\frac{3}{2}}}$。分子$a^2(1-\cos t)$，分母$a^3[2(1-\cos t)]^{\frac{3}{2}}$，所以$\kappa=\frac{1}{4a\sin(\frac{t}{2})}$。也可由$\kappa=\frac{|y_{xx}''|}{[1+(y_x')^2]^{\frac{3}{2}}}$核对。\par\textbf{方法点评} 参数公式在切线接近竖直时较稳健。\par\textbf{变式与易错点} 尖点$t=0,2\pi$速度为0，以上正规曲线公式不适用。\end{example}

\begin{theorem}{Leibniz公式}若$f,g$有$n$阶导数，则$(fg)^{(n)}=\sum_{k=0}^n\binom nkf^{(k)}g^{(n-k)}$。\end{theorem}\begin{proof}$n=1$是乘积法则。若$n$成立，再求导后将$f^{(k+1)}g^{(n-k)}$和$f^{(k)}g^{(n+1-k)}$按相同指标合并，用Pascal恒等式$\binom n{k-1}+\binom nk=\binom{n+1}k$得到$n+1$式。\end{proof}

\begin{example}{考研风格｜高阶导数的有限项化}{难度：$\star\star\star$}\label{deep:ex:differential_source_deep:leibniz}\textbf{题目} 求$(x^2\e^{ax})^{(n)}$，$n\ge2$。\par\textbf{思路} 多项式三阶以后导数为0。\par\textbf{详细解答} Leibniz公式只剩$k=0,1,2$三项：$\e^{ax}[a^nx^2+2na^{n-1}x+n(n-1)a^{n-2}]$。当$a=0$应直接由原函数区分$n=2$给2、$n>2$给0，以免书写$0^0$。另一法指数移位$D^n(\e^{ax}x^2)=\e^{ax}(D+a)^nx^2$，二项式展开同样只有三项。\par\textbf{方法点评} 两种表达本质相同，算子法方便与ODE章节衔接。\par\textbf{变式与易错点} 参数为0时从原式处理退化。\end{example}

\subsection{辅助函数：端点、零点与导数信号}\label{deep:differential_source_deep:mvt}

\begin{method}{辅助函数：端点、零点与导数信号}\textbf{适用条件} Rolle要求闭区间连续、开区间可微，端点值相等。\par\textbf{步骤与依据} 从目标导数反积分，或用指数因子消去$f'-\lambda f$；多点条件用插值多项式制造多个零点。\par\textbf{易错点} 只写“由中值定理”而不展示辅助函数，容易掩盖端点条件缺失。\end{method}

\begin{example}{竞赛风格｜指数辅助函数双思路}{难度：$\star\star\star$}\label{deep:ex:differential_source_deep:auxiliary}\textbf{题目} $f\in C[a,b]$且在$(a,b)$可微，$f(a)=f(b)=0$。证明对任意$\lambda\in\R$，存在$c$使$f'(c)=\lambda f(c)$。\par\textbf{思路} 目标是$(\e^{-\lambda x}f)'=0$。\par\textbf{详细解答} 令$g(x)=\e^{-\lambda x}f(x)$，端点均0，由Rolle存在$c$使$g'(c)=\e^{-\lambda c}[f'(c)-\lambda f(c)]=0$。另一思路：若$f'-\lambda f$处处不为零，则$g'=\e^{-\lambda x}(f'-\lambda f)$也处处不为零。由导数的Darboux性质，$g'$在$(a,b)$同号；再由中值定理知$g$严格单调，与$g(a)=g(b)=0$矛盾。此反证也不要求$f'$连续。\par\textbf{方法点评} 比较证明条件，可以避免为反证法无意增加假设。\par\textbf{变式与易错点} 不能在$f(c)=0$时除以$f(c)$写$\frac{f'}{f}=\lambda$。\end{example}

\begin{example}{竞赛风格｜插值余项与多点Rolle}{难度：$\star\star\star$}\label{deep:ex:differential_source_deep:interpolation}\textbf{题目} $f\in C^2[a,b],a<c<b$，证明存在$\xi\in(a,b)$使$f(c)-L(c)=\frac{f''(\xi)(c-a)(c-b)}{2}$，$L$为端点线性插值。\par\textbf{思路} 用二次因子制造三个零点。\par\textbf{详细解答} 设$K=\frac{f(c)-L(c)}{(c-a)(c-b)}$，$g(x)=f(x)-L(x)-K(x-a)(x-b)$。$g(a)=g(c)=g(b)=0$，两次Rolle给$g''(\xi)=0$。因$L''=0$、二次因子二阶导数为2，得$f''(\xi)=2K$，即结论。若$f''\ge0$，由$(c-a)(c-b)<0$知$f(c)\le L(c)$，也得到凸函数图像在弦下。\par\textbf{方法点评} 证明余项与凸性是同一个零点构造。\par\textbf{变式与易错点} $\xi$依赖$c$，不能把不同$c$对应的$\xi$视为同一固定点。\end{example}

\subsection{Taylor积分余项、凸性与不等式}\label{deep:differential_source_deep:taylor}

\begin{method}{Taylor积分余项、凸性与不等式}\textbf{适用条件} 导数具有足够连续性，或使用相应较弱余项版本。\par\textbf{步骤与依据} 将高阶导数积分回去，利用核符号给双边界；最值问题先检临界点再检端点。\par\textbf{易错点} 高阶导数在某点为0不自动说明拐点，须看符号变化。\end{method}

\begin{theorem}{Taylor积分余项}若$f\in C^{n+1}$，则$f(x)=\sum_{k=0}^n\frac{f^{(k)}(a)(x-a)^k}{k!}+\frac1{n!}\int_a^x(x-t)^nf^{(n+1)}(t)\dd t$。\end{theorem}\begin{proof}$n=0$是基本定理。对余项分部一次，将$f^{(n+1)}$降为$f^{(n)}$，边界产生下一Taylor项；重复$n$次即还原$f(x)-f(a)$。反向也可对右边逐次求导，并比较在$a$的初值。\end{proof}

\begin{example}{考研风格｜对数不等式双解}{难度：$\star\star$}\label{deep:ex:differential_source_deep:logbound}\textbf{题目} 证明$x>0$时$\frac{x}{1+x}<\ln(1+x)<x$。\par\textbf{思路} 用积分夹逼或导数单调。\par\textbf{详细解答} 积分法：在$0<t<x$，$\frac{1}{1+x}<\frac{1}{1+t}<1$，积分得结论。导数法：$g=x-\ln(1+x)$有$g(0)=0,g'=\frac{x}{1+x}>0$；$h=\ln(1+x)-\frac{x}{1+x}$有$h(0)=0,h'=\frac{x}{(1+x)^2}>0$。\par\textbf{方法点评} 积分法最短，导数法适合构造更高阶余量。\par\textbf{变式与易错点} $-1<x<0$时积分方向反转，不能原样沿用同一链条。\end{example}

\begin{example}{考研风格｜渐近线与精细差值}{难度：$\star\star\star$}\label{deep:ex:differential_source_deep:asymptote}\textbf{题目} 求$f(x)=\sqrt{x^2+x+1}$在$\pm\infty$的斜渐近线，并给首项误差。\par\textbf{思路} 在两端分别提取$|x|$。\par\textbf{详细解答} $x\to+\infty$，$f=x\sqrt{1+\frac{1}{x}+\frac{1}{x^2}}=x+\frac{1}{2}+\frac{3}{8x}+O(x^{-2})$，故渐近线$y=x+\frac{1}{2}$。$x\to-\infty$，$f=-x\sqrt{1+\frac{1}{x}+\frac{1}{x^2}}=-x-\frac{1}{2}-\frac{3}{8x}+O(x^{-2})$，故$y=-x-\frac{1}{2}$。也可先求斜率极限$\frac{f}{x}=\pm1$，再分别有理化求截距。\par\textbf{方法点评} 两端斜率不同源于$|x|$。\par\textbf{变式与易错点} 把$\sqrt{x^2}$写成$x$会丢失负无穷方向。\end{example}

\section{图解：割线趋向切线}
\begin{figure}[htbp]\centering
\begin{tikzpicture}\begin{axis}[width=.8\textwidth,height=6cm,axis lines=middle,xmin=0,xmax=2.2,ymin=-.5,ymax=4.2,xlabel=$x$,ylabel=$y$,samples=100,legend pos=north west]
\addplot[Teal,thick,domain=0:2]{x^2};\addlegendentry{$y=x^2$}
\addplot[Ink,dashed,domain=.4:2]{2*x-1};\addlegendentry{$y=2x-1$}
\addplot[Rust,domain=.65:2]{2.8*x-1.8};\addlegendentry{割线}
\fill[Ink] (axis cs:1,1) circle[radius=1.8pt];\fill[Ink] (axis cs:1.8,3.24) circle[radius=1.8pt];
\draw[<->] (axis cs:1,.5)--node[below]{$h$}(axis cs:1.8,.5);
\end{axis}\end{tikzpicture}
\caption{割线斜率$\frac{(1+h)^2-1}{h}=2+h$，当$h\to0$趋向切线斜率2。切线是局部线性近似，实际误差$(1+h)^2-(1+2h)=h^2$。图取$h=0.8$以便看清两点。}\label{fig:review:derivative}
\end{figure}

\section{微分中的反推、分割与几何辨析}\label{cd:differential}
\begin{examtip}
本节集中训练容易被一般公式掩盖的结构：用伸缩差商反推可导性、按函数值分割区间、构造自动抵消的辅助函数、控制中值点位置，以及区分一点导数信息与整个邻域的几何性质。
\end{examtip}

\subsection{非对称伸缩差商反推可导}
\begin{method}{伸缩差商的几何迭代}
\textbf{适用条件} 已知$f$在$a$连续，且$f(a+\alpha t)-f(a+\beta t)$的差商有有限极限，$|\alpha|\ne|\beta|$。\par
\textbf{步骤} 先将较大尺度记为$h$，把较小尺度写成$rh$，其中$|r|<1$；反复缩放并望远镜求和；用连续性消去最终的函数增量。\par
\textbf{疑难点} 对称尺度$\beta=-\alpha$只检测奇部，可能完全漏掉尖点。
\end{method}
\begin{theorem}{非等模尺度的可导判据}\label{cd:dilation-test}
设$f$在$a$的邻域内有定义并在$a$连续，$\alpha,\beta\in\R$且$|\alpha|\ne|\beta|$。若
\[
 \lim_{t\to0}\dfrac{f(a+\alpha t)-f(a+\beta t)}{t}=L,
\]
则$f'(a)$存在，且$f'(a)=\frac{L}{\alpha-\beta}$。
\end{theorem}
\begin{proof}
交换$\alpha,\beta$只使$L$变号，可设$|\alpha|>|\beta|$。记$r=\frac{\beta}{\alpha}$，$F(h)=f(a+h)-f(a)$，则$|r|<1$。若$r=0$，结论就是导数定义。以下设$r\ne0$，令
\[
 D(h)=\dfrac{F(h)-F(rh)}{h},\qquad D(h)\to\ell=\dfrac L\alpha.
\]
对固定充分小的$h\ne0$，望远镜恒等式给
\[
 \dfrac{F(h)}h
 =\sum_{k=0}^{N-1}r^kD(r^kh)+\dfrac{F(r^Nh)}h.
\]
由$f$在$a$连续，最后一项随$N\to\infty$趋0。给定$\varepsilon>0$，取$\delta>0$使$0<|u|<\delta$时$|D(u)-\ell|<\varepsilon$。若$0<|h|<\delta$，所有$r^kh$都在该范围，故级数绝对收敛，而且
\[
 \left|\dfrac{F(h)}h-\dfrac{\ell}{1-r}\right|
 \le\sum_{k=0}^{\infty}|r|^k\varepsilon
 =\dfrac{\varepsilon}{1-|r|}.
\]
令$\varepsilon\downarrow0$得到导数极限，$\frac{\ell}{1-r}=\frac{L}{\alpha-\beta}$。
\end{proof}
\begin{example}{竞赛风格｜两个尺度能否确定导数}{难度：$\star\star\star$}\label{cd:ex:dilation}
\textbf{题目} 已知$f$在0连续，且$\lim_{h\to0}\frac{f(3h)-f(-h)}{h}=8$，求$f'(0)$。若只知$\lim_{h\to0}\frac{f(h)-f(-h)}{2h}$存在，结论是否仍成立？\par
\textbf{思路} 第一问尺度模不同，第二问尺度模相同。\par
\textbf{详细解答} 由定理\ref{cd:dilation-test}，第一问$f'(0)=\frac{8}{3-(-1)}=2$。也可直接取$r=-\frac{1}{3}$，将$f(h)-f(0)$拆为几何尺度上的连续差，得到极限2。第二问取$f(x)=|x|$，则每个对称差商都为0，但右导数为1、左导数为$-1$，所以$f'(0)$不存在。\par
\textbf{方法点评} 第一问并未预先知道$f$可导，因此不能先把两个增量各自替换成$f'(0)$乘增量；缩放迭代正是反推可导性的关键。\par
\textbf{疑难点与变式} 若取消$f$在0连续，改变单独一个点值不影响伸缩差在充分小非零$h$处的值，却能破坏可导性。
\end{example}

\subsection{零点因子对尖点的消除}
\begin{method}{判断乘积是否消除不可导}
\textbf{适用条件} $u$在$a$可导，$v$在$a$连续而不可导。\par
\textbf{步骤} 若$u(a)=0$，直接用差商证明$uv$可导；若$u(a)\ne0$，用商$\frac{uv}{u}$反证。\par
\textbf{疑难点} “连续而不可导”与“无界振荡”不同；没有连续性，乘以一个零点未必能保证所需可导阶数。
\end{method}
\begin{theorem}{连续尖点的乘积判据}
在上述条件下，$uv$在$a$可导当且仅当$u(a)=0$；此时$(uv)'(a)=u'(a)v(a)$。
\end{theorem}
\begin{proof}
若$u(a)=0$，则
\[
 \dfrac{u(a+h)v(a+h)-u(a)v(a)}h
 =\dfrac{u(a+h)-u(a)}h\,v(a+h)\to u'(a)v(a).
\]
若$u(a)\ne0$而$uv$可导，由连续性$u$在$a$附近非零，可导函数之商$v=\frac{uv}{u}$在$a$可导，与假设矛盾。
\end{proof}
\begin{example}{考研风格｜多项式因子消去部分尖点}{难度：$\star\star$}\label{cd:ex:kinks}
\textbf{题目} 求$F(x)=(x^2+x-2)|x^3-x|$的全部不可导点。\par
\textbf{思路} 绝对值内部非零时光滑；只需检查$x=-1,0,1$。\par
\textbf{详细解答} $p(x)=x^3-x$在这三点都是简单零点，因为$p'(\pm1)=2$、$p'(0)=-1$。因此$|p|$在每个零点连续而不可导，其左右导数为一对相反的非零数。另一方面
\[
 x^2+x-2=(x+2)(x-1)
\]
只在$x=1$同时消去尖点。由乘积判据，$F$在1可导，且$F'(1)=0$；在$-1,0$不可导。其余各点绝对值符号在局部固定，$F$光滑。\par
\textbf{方法点评} 先列绝对值的简单零点，再检查外因子是否同时为零，远比展开四个分段多项式简洁。\par
\textbf{疑难点与变式} 若绝对值内部有偶重零点，绝对值本身可能光滑，应先看局部零点重数，不能把所有根都列成尖点。
\end{example}

\subsection{按函数值分割产生多个中值点}
\begin{method}{先规定增量，再寻找分割点}
\textbf{适用条件} 要构造多个不同点的导数关系，已知$f(0)=0,f(1)=1$且$f$连续可导。\par
\textbf{步骤} 根据目标权重规定一串严格递增的函数值；依次用介值定理选出有序分割点；在相邻小区间分别用Lagrange中值定理。\par
\textbf{疑难点} 不能随意选横坐标再要求函数值按比例变化；本方法先分函数值，并不要求$f$单调。
\end{method}
\begin{example}{竞赛风格｜倒数导数的加权和}{难度：$\star\star\star$}\label{cd:ex:weighted-reciprocal}
\textbf{题目} 设$f\in C[0,1]$且在$(0,1)$可导，$f(0)=0,f(1)=1$。给定正数$w_1,\ldots,w_m$，证明存在互异$\xi_1,\ldots,\xi_m\in(0,1)$，且$f'(\xi_j)>0$，满足
\[
 \sum_{j=1}^m\dfrac{w_j}{f'(\xi_j)}
 =\sum_{j=1}^mw_j.
\]
\textbf{思路} 让第$j$段的函数增量恰等于权重在总权重中的比例。\par
\textbf{详细解答} 记$W=\sum w_j$，$c_j=\frac{w_1+\cdots+w_j}{W}$，$c_0=0,c_m=1$。从$t_0=0$开始，对$j=1,\ldots,m-1$依次选点：若已选$t_{j-1}<1$且$f(t_{j-1})=c_{j-1}$，由于$c_{j-1}<c_j<1=f(1)$，介值定理在$(t_{j-1},1)$中给出$t_j$使$f(t_j)=c_j$；最后置$t_m=1$。于是
\[
 0=t_0<t_1<\cdots<t_m=1.
\]
对每段应用Lagrange定理，存在$\xi_j\in(t_{j-1},t_j)$使
\[
 f'(\xi_j)=\dfrac{f(t_j)-f(t_{j-1})}{t_j-t_{j-1}}
 =\dfrac{w_j}{W(t_j-t_{j-1})}>0.
\]
所以$\frac{w_j}{f'(\xi_j)}=W(t_j-t_{j-1})$。求和后区间长度望远镜相消，结果为$W$。各$\xi_j$位于不相交的开区间，故互异。\par
\textbf{方法点评} 中值点的互异性来自分段，而导数非零来自严格增加的目标函数值；这两项都应在除法前说明。\par
\textbf{疑难点与变式} $w_j=1$时得到$\sum\frac{1}{f'(\xi_j)}=m$。若$f$振荡，同一函数值可能有多个逆像，只需按上述顺序逐次选取，不需要全局反函数。
\end{example}

\subsection{利用变上限积分自动抵消}
\begin{method}{让乘积求导的中间项相消}
\textbf{适用条件} 目标包含$f'(\xi)\int_\xi^bg$与$g'(\xi)\int_a^\xi f$。\par
\textbf{步骤} 同时构造$f(x)\int_x^bg$和$g(x)\int_a^xf$；利用积分上下限求导的相反符号消掉$fg$；由边界条件制造两个零点。\par
\textbf{疑难点} 下限为$x$时导数带负号；遗漏该负号会使两项相加而无法抵消。
\end{method}
\begin{example}{竞赛风格｜交叉积分的Rolle构造}{难度：$\star\star\star$}\label{cd:ex:cross-integral}
\textbf{题目} 设$f,g\in C[a,b]$，在$(a,b)$可导，且$f(a)=g(b)=0$。证明存在$\xi\in(a,b)$，使
\[
 f'(\xi)\int_\xi^bg(t)\dd t
 +g'(\xi)\int_a^\xi f(t)\dd t=0.
\]
\textbf{思路} 将目标看成一个乘积和的导数，补出的两项应恰好相消。\par
\textbf{详细解答} 定义
\[
 F(x)=f(x)\int_x^bg(t)\dd t
      +g(x)\int_a^xf(t)\dd t.
\]
$F$在$[a,b]$连续、在$(a,b)$可导，$F(a)=F(b)=0$。由乘积法则和微积分基本定理，
\begin{align*}
 F'(x)
 &=f'(x)\int_x^bg(t)\dd t-f(x)g(x)\\
 &\quad+g'(x)\int_a^xf(t)\dd t+g(x)f(x)\\
 &=f'(x)\int_x^bg(t)\dd t+g'(x)\int_a^xf(t)\dd t.
\end{align*}
Rolle定理给$F'(\xi)=0$，即结论。\par
\textbf{方法点评} 构造的目标不是把所有项分别积分，而是使求导产生的额外项互相抵消。\par
\textbf{疑难点与变式} 没有假设$f',g'$连续，直接对$F$使用Rolle即可；不应额外假设目标表达式连续后再用介值定理。
\end{example}

\subsection{能量辅助函数与退化因子的排除}
\begin{method}{先迫使辅助函数取内点极值}
\textbf{适用条件} 目标含$f+f''$，且已知$f$的全局界及某点$f^2+(f')^2$较大。\par
\textbf{步骤} 构造$E=f^2+(f')^2$；用两侧中值定理找出$E$较小的端点；对内点最大值求导；用最大值大小排除$f'=0$。\par
\textbf{疑难点} 从$E'=2f'(f+f'')=0$不能直接删除$f'$，必须证明所取极值点处$f'$非零。
\end{method}
\begin{example}{竞赛风格｜有界振幅迫使二阶组合为零}{难度：$\star\star\star$}\label{cd:ex:energy}
\textbf{题目} 设$f\in C^2[-2,2]$，$|f(x)|\le1$，且$f(0)^2+f'(0)^2=4$。证明某$\xi\in(-2,2)$满足$f(\xi)+f''(\xi)=0$。\par
\textbf{思路} 在0两侧先找斜率较小的点，再比较能量的最大值。\par
\textbf{详细解答} 在$[-2,0]$应用Lagrange定理，存在$u\in(-2,0)$使
\[
 |f'(u)|=\dfrac{|f(0)-f(-2)|}{2}\le1.
\]
同理存在$v\in(0,2)$使$|f'(v)|\le1$。令$E(x)=f(x)^2+f'(x)^2$，则$E(u),E(v)\le2$，而$E(0)=4$。连续函数$E$在$[u,v]$取得最大值，其最大点$\xi$必在$(u,v)$。Fermat定理给
\[
 0=E'(\xi)=2f'(\xi)\bigl(f(\xi)+f''(\xi)\bigr).
\]
若$f'(\xi)=0$，则$E(\xi)=f(\xi)^2\le1$，与$E(\xi)\ge E(0)=4$矛盾。因此$f'(\xi)\ne0$，从而$f(\xi)+f''(\xi)=0$。\par
\textbf{方法点评} 数值4同时承担“能量大于边界值”和“排除退化因子”的作用。\par
\textbf{疑难点与变式} 辅助函数有极值不够，必须确保极值在内部。若只知$E(0)=1$，排除$f'(\xi)=0$的论证便不能照用。
\end{example}

\subsection{参数竖切线处的凹凸改变}
\begin{method}{参数导数失效点返回几何定义}
\textbf{适用条件} 参数曲线$\bm r(t)$正规，但$x'(t_0)=0$。\par
\textbf{步骤} 检查$y'(t_0)$判断竖切线；核对$x(t)$在附近是否严格单调；在$t_0$两侧分别计算$y_{xx}$并比较凹凸，必要时用切线两侧的位置核对。\par
\textbf{疑难点} $x'=0$使$y_x=\frac{y'}{x'}$失效，却不意味着曲线一定是尖点；正规性要求$(x',y')\ne(0,0)$。
\end{method}
\begin{example}{考研风格｜二阶导数不存在的凹凸转折}{难度：$\star\star\star$}\label{cd:ex:vertical-inflection}
\textbf{题目} 曲线由$x=t^3,\ y=t+t^2$给出。分析$t=0$处切线及两侧凹凸。\par
\textbf{思路} $x(t)$严格递增，所以$t$两侧与$x$两侧对应；在非零$t$处用参数求导。\par
\textbf{详细解答} $x'=3t^2,y'=1+2t$，在$t=0$速度向量为$(0,1)$，故曲线正规，切线为$x=0$。当$t\ne0$，
\[
 y_x=\dfrac{1+2t}{3t^2},\qquad
 y_{xx}=\dfrac{-2(1+t)}{9t^5}.
\]
取$0<|t|<\frac{1}{2}$，则$1+t>0$：左侧$t<0$时$y_{xx}>0$，右侧$t>0$时$y_{xx}<0$。因此两侧凹凸改变，且曲线穿过竖切线。若将有竖切线且两侧凹凸改变的点纳入几何拐点定义，则原点为拐点。\par
\textbf{方法点评} 凹凸判别在两侧分别进行，不要求二阶导数在转折点存在。这里显式函数为$y=x^{\frac{1}{3}}+x^{\frac{2}{3}}$，可独立核对两侧符号。\par
\textbf{疑难点与变式} 有些课程只将具有有限切线斜率的转折点称为拐点。遇到这种约定，应报告“原点具有竖切线且两侧凹凸改变”，避免混淆定义。
\end{example}
\begin{center}
\begin{tikzpicture}[x=3.0cm,y=1.8cm]
 \draw[->,gray] (-1.15,0)--(1.2,0) node[right] {$x$};
 \draw[->,gray] (0,-.45)--(0,2.35) node[above] {$y$};
 \draw[very thick,Teal,domain=-1:1,samples=100,variable=\t]
   plot ({\t*\t*\t},{\t+\t*\t});
 \draw[dashed,Rust] (0,-.4)--(0,2.2);
 \fill[Ink] (0,0) circle (1.5pt);
 \node[below right] at (0,0) {$O$};
 \node[align=left,font=\small] at (-.65,.65) {$t<0$\\$y_{xx}>0$};
 \node[align=left,font=\small] at (.65,1.25) {$t>0$\\$y_{xx}<0$};
 \node[align=left,font=\small,anchor=west] at (1.3,1.6)
   {切线为$x=0$\\速度$(0,1)\ne(0,0)$\\原点处曲线正规};
\end{tikzpicture}
\par\small 图：参数$t$从左向右增加。虚线是原点处的竖切线；左右两支的凹凸符号相反。图示用于区分竖切线与速度为零的尖点。
\end{center}

\subsection{Taylor展开中值点本身}
\begin{method}{把未知中值点写成相对位置}
\textbf{适用条件} $f\in C^n$，在$a$处第2至$n-1$阶导数为0，第$n$阶导数非零。\par
\textbf{步骤} 写中值点为$a+\theta_hh$，$0<\theta_h<1$；分别展开$f(a+h)$和$f'(a+\theta_hh)$；比较首个非零的高阶项。\par
\textbf{疑难点} $\theta_h$随$h$变化，不能先假设它收敛。应先由等式推出$\theta_h^{n-1}$的极限。
\end{method}
\begin{theorem}{Lagrange中值点的高阶定位}\label{cd:mvt-position}
设$n\ge2$，$f$在$a$附近属于$C^n$，$f^{(k)}(a)=0$对$2\le k\le n-1$成立，且$f^{(n)}(a)\ne0$。对每个充分小的$h\ne0$，任取满足
\[
 f(a+h)-f(a)=h f'(a+\theta_hh),\qquad0<\theta_h<1
\]
的中值点，则$\theta_h\to n^{-\frac{1}{n-1}}$。
\end{theorem}
\begin{proof}
记$c=f^{(n)}(a)\ne0$。Taylor展开得
\[
 f(a+h)-f(a)=f'(a)h+\dfrac{c}{n!}h^n+o(h^n),
\]
而
\[
 f'(a+\theta_hh)
 =f'(a)+\dfrac{c}{(n-1)!}(\theta_hh)^{n-1}+o(h^{n-1}).
\]
第二式的余项可统一写成$o(h^{n-1})$：将点余项写为$\eta(s)s^{n-1}$，其中$\eta(s)\to0$，再用$|s|=|\theta_hh|\le|h|$即可。代入中值等式并除以$h^n$，得
\[
 \dfrac{c}{n!}+o(1)
 =\dfrac{c}{(n-1)!}\theta_h^{\,n-1}+o(1).
\]
故$\theta_h^{n-1}\to\frac{1}{n}$。因$\theta_h>0$，取正的$(n-1)$次根即得结论。
\end{proof}
\begin{example}{竞赛风格｜反正切中值点的位置}{难度：$\star\star\star$}\label{cd:ex:mvt-position}
\textbf{题目} 对$x\ne0$，设$\xi_x$介于0与$x$之间且$\arctan x=\frac{x}{1+\xi_x^2}$。求$\lim_{x\to0}\frac{\xi_x}{x}$。\par
\textbf{思路} 可用高阶定位定理，也可从等式直接解出平方比。\par
\textbf{详细解答} $f(x)=\arctan x$满足$f'(0)=1,f''(0)=0,f'''(0)=-2\ne0$。由定理\ref{cd:mvt-position}，取$n=3$，得$\frac{\xi_x}{x}\to\frac{1}{\sqrt3}$。独立计算：题设等式给
\[
 \dfrac{\xi_x^2}{x^2}
 =\dfrac{x-\arctan x}{x^2\arctan x}
 \to\dfrac13,
\]
其中$\arctan x=x-\frac{x^3}{3}+o(x^3)$。由于$\xi_x$与$x$同号，$\frac{\xi_x}{x}>0$，故仍得$\frac{1}{\sqrt3}$。\par
\textbf{方法点评} 平方比的极限只确定绝对值，正号来自中值点所在区间。\par
\textbf{疑难点与变式} 不能仅从$|\xi_x|<|x|$推出$\frac{\xi_x}{x}$有极限；夹在0和1之间只表示有界。
\end{example}

\subsection{用所有小区间的积分平均反推凸性}
\begin{method}{对称小区间检测二阶项}
\textbf{适用条件} $f\in C^2(I)$，给出对每个子区间成立的积分平均不等式。\par
\textbf{步骤} 取以任意$c$为中心的小区间$[c-h,c+h]$；展开后积分，奇次项抵消；二阶项系数的符号给出$f''(c)$。\par
\textbf{疑难点} 单独一个区间的不等式不能控制所有点的凹凸性；“每个子区间”是反推的关键。
\end{method}
\begin{theorem}{中点积分平均的凸性判据}\label{cd:mean-convex}
设$I$为开区间、$f\in C^2(I)$。以下两条件等价：
\[
 f''(x)\ge0\quad(x\in I);
 \qquad
 f\!\left(\dfrac{a+b}{2}\right)
 \le\dfrac1{b-a}\int_a^bf(x)\dd x
 \quad([a,b]\subset I).
\]
\end{theorem}
\begin{proof}
若$f''\ge0$，则$f$在每点$c$的切线位于图像下方，
$f(x)\ge f(c)+f'(c)(x-c)$。取$c=\frac{a+b}{2}$并积分，线性项积分为0，得到右侧条件。反向固定$c\in I$，由Taylor公式，
\[
 f(c+t)=f(c)+f'(c)t+\dfrac{f''(c)}2t^2+r(t),
 \qquad r(t)=o(t^2).
\]
对$[-h,h]$积分并除以$2h$，余项满足
$\left|(2h)^{-1}\int_{-h}^hr(t)\dd t\right|=o(h^2)$，所以
\[
 \dfrac1{2h}\int_{c-h}^{c+h}f(x)\dd x-f(c)
 =\dfrac{f''(c)}6h^2+o(h^2).
\]
左边非负，除以$h^2$后令$h\to0^+$得$f''(c)\ge0$。$c$任意，结论成立。
\end{proof}
\begin{example}{竞赛风格｜积分平均判据的必要范围}{难度：$\star\star\star$}\label{cd:ex:mean-convex}
\textbf{题目} 设$f(x)=x^4-\lambda x^2$。求使定理\ref{cd:mean-convex}中的积分平均不等式在每个实数子区间上成立的全部$\lambda$。若只在$[-1,1]$检查，条件有何变化？\par
\textbf{思路} 全部区间用二阶导数，单个区间直接积分。\par
\textbf{详细解答} $f''(x)=12x^2-2\lambda$，对所有$x$非负当且仅当$\lambda\le0$。因此全部子区间的条件是$\lambda\le0$。若只检查$[-1,1]$，中点值$f(0)=0$，而平均值为
\[
 \dfrac12\int_{-1}^1(x^4-\lambda x^2)\dd x
 =\dfrac15-\dfrac{\lambda}{3}.
\]
只须$\lambda\le\frac{3}{5}$；例如$\lambda=\frac{1}{2}$通过这一个区间的检验，却有$f''(0)=-1<0$。\par
\textbf{方法点评} 局部二阶项可由任意小区间的平均提取，单个大区间则可能由其他位置的正贡献掩盖局部负曲率。\par
\textbf{疑难点与变式} “积分非负”一般不推出被积函数逐点非负；只有对足够丰富的小区间都成立，才可利用连续性反推逐点信息。
\end{example}

\subsection{指数因子化简截断多项式的零点}
\begin{method}{消去相邻阶数，只留下最高次项}
\textbf{适用条件} $P_n(x)=\sum_{k=0}^n\frac{x^k}{k!}$，有$P_n'=P_{n-1}$。\par
\textbf{步骤} 构造$\e^{-x}P_n$，求导后只剩$-\frac{\e^{-x}x^n}{n!}$；按$n$的奇偶分区间判断单调和端点极限。\par
\textbf{疑难点} $P_n$本身的导数符号未必容易判断，乘以处处为正的指数因子既不改变零点，又使求导大幅简化。
\end{method}
\begin{example}{竞赛风格｜指数部分和的实根个数}{难度：$\star\star\star$}\label{cd:ex:exp-polynomial}
\textbf{题目} 对整数$n\ge1$，判断$P_n(x)=1+x+\cdots+\frac{x^n}{n!}$的实根个数。\par
\textbf{思路} 保留零点，改变用于判断单调的函数。\par
\textbf{详细解答} 令$F_n(x)=\e^{-x}P_n(x)$。因$\e^{-x}>0$，两者实根完全相同，且
\[
 F_n'(x)=\e^{-x}(P_n'-P_n)
 =-\dfrac{x^n}{n!}\e^{-x}.
\]
若$n$为偶数，$F_n$在$\R$严格递减，$\lim_{x\to+\infty}F_n(x)=0$，故每个有限$x$处$F_n(x)>0$，没有实根。若$n$为奇数，则$F_n$在$(-\infty,0)$严格递增、在$(0,\infty)$严格递减，$F_n(0)=1$。当$x\to-\infty$时，$P_n$的最高次项为负而$\e^{-x}\to+\infty$，故$F_n(x)\to-\infty$。介值定理及严格单调性给负半轴恰有一个根。正半轴上$P_n$的每项均正，所以没有根。\par
\textbf{方法点评} 根的存在性由两端符号给出，唯一性由变换后的单调性给出；不必求解高次多项式。\par
\textbf{疑难点与变式} 偶数情形“递减且趋0”给的是严格正值，不能把极限0算成一个有限实根。
\end{example}

\subsection{一点的二阶导数不能替代邻域条件}
\begin{method}{反例中分开控制低阶极限与高阶振荡}
\textbf{适用条件} 需要辨析$f''(a)>0$与邻域内$f''>0$的差异。\par
\textbf{步骤} 以$x^2$固定原点的二阶导数，再叠加高次小量与快速振荡；计算一点的导数时用差商，在非零点才套求导法则。\par
\textbf{疑难点} 在一点$f''(a)>0$足以配合$f'(a)=0$判严格极小，却不足以保证整个邻域凸。
\end{method}
\begin{example}{考研风格｜有严格极小却没有凸邻域}{难度：$\star\star\star$}\label{cd:ex:point-convex}
\textbf{题目} 令$f(0)=0$，$x\ne0$时$f(x)=x^2+x^5\sin(\frac{1}{x^2})$。证明$f''(0)=2$、原点为严格极小点，但$f$在原点任一邻域内都不是凸函数。\par
\textbf{思路} 极小值用函数值估计，二阶导数的连续性失败用振荡子列展示。\par
\textbf{详细解答} 由定义，$f'(0)=\lim_{h\to0}[h+h^4\sin(\frac{1}{h^2})]=0$。对$x\ne0$，
\[
 f'(x)=2x+5x^4\sin(\dfrac{1}{x^2})-2x^2\cos(\dfrac{1}{x^2}).
\]
故$f''(0)=\lim_{h\to0}\frac{f'(h)}{h}=2$。同时$0<|x|<1$时
\[
 f(x)\ge x^2-|x|^5=x^2(1-|x|^3)>0=f(0),
\]
所以0是严格极小点。在非零点继续求导，
\[
 f''(x)=2+20x^3\sin(\dfrac{1}{x^2})
          -14x\cos(\dfrac{1}{x^2})-\dfrac4x\sin(\dfrac{1}{x^2}).
\]
取$x_k=(\frac{\pi}{2}+2k\pi)^{-\frac{1}{2}}\to0^+$，则$f''(x_k)=2+20x_k^3-\frac{4}{x_k}\to-\infty$。任意原点邻域内都出现$f''<0$的点，因二阶可导凸函数必满足$f''\ge0$，故不存在凸邻域。\par
\textbf{方法点评} 原点的低阶差商与邻近点的高阶振荡可以同时成立，不应把一点导数值自动推广为整个邻域的导数符号。\par
\textbf{疑难点与变式} 若再假设$f''$在0连续，则$f''(0)>0$可由保号性推出一个邻域内$f''>0$；连续性正是缺少的桥梁。
\end{example}

\subsection{从端点水平切线推出二阶变化下界}
\begin{method}{同一中点作两端Taylor展开}
\textbf{适用条件} $f\in C^2[a,b]$且$f'(a)=f'(b)=0$。\par
\textbf{步骤} 在$a,b$分别向同一个中点展开；相减消去中点函数值；把两处二阶导数绝对值之和换为其中较大者的两倍。\par
\textbf{疑难点} 两次展开的余项点不同；可以比较它们的最大值，却不能将两点直接写为同一点。
\end{method}
\begin{example}{竞赛风格｜端点平坦仍须有足够的弯曲}{难度：$\star\star\star$}\label{cd:ex:endpoint-bound}
\textbf{题目} 设$f\in C^2[a,b]$，$a<b$，$f'(a)=f'(b)=0$。证明存在$\xi\in(a,b)$，使
\[
 |f''(\xi)|\ge\dfrac{4|f(b)-f(a)|}{(b-a)^2}.
\]
\textbf{思路} 中点到两端距离相等，使两个二阶余项拥有相同系数。\par
\textbf{详细解答} 记$m=\frac{a+b}{2}$，$L=b-a$。Taylor公式分别给出$u\in(a,m)$、$v\in(m,b)$，使
\[
 f(m)=f(a)+\dfrac{L^2}{8}f''(u),\qquad
 f(m)=f(b)+\dfrac{L^2}{8}f''(v).
\]
相减得$f(b)-f(a)=\frac{L^2[f''(u)-f''(v)]}{8}$。因此
\[
 |f(b)-f(a)|
 \le\dfrac{L^2}{8}\bigl(|f''(u)|+|f''(v)|\bigr)
 \le\dfrac{L^2}{4}\max\{|f''(u)|,|f''(v)|\}.
\]
从$u,v$中取绝对值较大的一点作为$\xi$即可。\par
\textbf{方法点评} 与先在整个区间只用一次Taylor相比，两端向中点展开把距离平方减小到$\frac{L^2}{4}$，得到更强系数。\par
\textbf{疑难点与变式} 若$f(a)=f(b)$，右端为0，结论仍成立；若取消端点导数为0，线性函数就能使右端正而$f''$恒为0。
\end{example}

\section{数三应用：边际量、需求弹性与收益}\label{r2:economics}
设产量为 $q$，总成本 $C(q)$、总收益 $R(q)$，总利润 $\Pi(q)=R(q)-C(q)$。边际成本、边际收益、边际利润分别是 $C'(q),R'(q),\Pi'(q)$，单位为“货币/单位产量”。对小增量 $\Delta q$，$\Delta C\approx C'(q)\Delta q$，不是把导数当成实际增加一单位的精确差额。

平均成本 $A(q)=C(q)/q$（$q>0$）满足
\[
 A'(q)=\frac{qC'(q)-C(q)}{q^2}=\frac{C'(q)-A(q)}q.
\]
因此平均成本内点驻点满足边际成本等于平均成本。利润内点候选满足 $R'=C'$，仍须检查二阶或导数变号，并与可行区间的边界比较。

\begin{example}{从总成本到平均成本与利润最大化}{数三应用}
令 $C(q)=100+10q+q^2$。则边际成本为 $10+2q$，平均成本为 $100/q+10+q$，其导数 $1-100/q^2$ 在 $q=10$ 两侧由负变正，故最小平均成本为 $30$，恰等于此时边际成本。
若价格固定为 $50$，则 $R=50q$，利润 $\Pi=40q-q^2-100$，$q\ge0$。$\Pi'=40-2q$，$\Pi''=-2$，所以最优产量为 $20$，最大利润为 $300$。平均成本最低的产量和利润最大的产量不是同一个问题。
\end{example}

\subsection{需求价格弹性：先说清符号约定}
设需求 $q=D(p)>0$、价格 $p>0$。有符号价格弹性定义为
\[
 E_{q,p}=\frac{pD'(p)}{D(p)}.
\]
对递减需求，它通常为负；不少经济学教材另用正值弹性 $\eta=-E_{q,p}$。收益 $R(p)=pD(p)$ 的导数与弹性为
\[
 R'(p)=D(p)+pD'(p)=D(p)(1-\eta),\qquad
 E_{R,p}=\frac{pR'(p)}{R(p)}=1-\eta.
\]
$\eta<1$ 时小幅提价增加收益，$\eta>1$ 时小幅提价减少收益；$\eta=1$ 仅给出内点驻点，还须判断是否最大。收益对价格的弹性与需求对价格的弹性不可混写。

\begin{example}{线性需求的两种变量写法}{数三应用}
令 $D(p)=120-2p$，$0<p<60$。则 $\eta=2p/(120-2p)$，收益 $R=120p-2p^2$ 在 $p=30$ 最大，此时需求 $60$、弹性为一。若反写 $p=60-q/2$，则 $R(q)=60q-q^2/2$，边际收益为 $60-q$，仍在 $q=60$ 取最大。
如果弹性用需求作为价格的函数定义，则小比例变化满足 $\Delta q/q\approx E_{q,p}\Delta p/p$；不能漏掉负号，也不能把百分比变化当成绝对数量变化。
\end{example}

\paragraph{自测与解答}
令 $D(p)=kp^{-a}$，$k,a>0$。则 $\eta=a$ 恒定，收益 $kp^{1-a}$ 在 $a<1$ 时递增、$a>1$ 时递减、$a=1$ 时恒定。单位弹性可能对应恒定收益，不能单靠 $\eta=1$ 宣称存在严格极大点。


\section*{章末分层习题}\addcontentsline{toc}{section}{章末分层习题}
\refstepcounter{exercise}\label{q:2:1}\par\noindent\textbf{习题\theexercise\quad 基础题} 求$\ln(1+x^2)$的导数。\quad\hyperref[a:2:1]{[答案]}\par
\refstepcounter{exercise}\label{q:2:2}\par\noindent\textbf{习题\theexercise\quad 基础题} 求$x\e^{2x}$的一阶导数。\quad\hyperref[a:2:2]{[答案]}\par
\refstepcounter{exercise}\label{q:2:3}\par\noindent\textbf{习题\theexercise\quad 基础题} 判断$|x|$在0是否可导。\quad\hyperref[a:2:3]{[答案]}\par
\refstepcounter{exercise}\label{q:2:4}\par\noindent\textbf{习题\theexercise\quad 考研风格} 求$(x\e^x)^{(n)}$，$n\ge1$。\quad\hyperref[a:2:4]{[答案]}\par
\refstepcounter{exercise}\label{q:2:5}\par\noindent\textbf{习题\theexercise\quad 考研风格} 求$x+\frac{1}{x}$在$(0,\infty)$上的最小值。\quad\hyperref[a:2:5]{[答案]}\par
\refstepcounter{exercise}\label{q:2:6}\par\noindent\textbf{习题\theexercise\quad 考研风格} 求$y=\ln x$在$x=1$处曲率。\quad\hyperref[a:2:6]{[答案]}\par
\refstepcounter{exercise}\label{q:2:7}\par\noindent\textbf{习题\theexercise\quad 竞赛拓展} 证明$|\sin x-\sin y|\le|x-y|$。\quad\hyperref[a:2:7]{[答案]}\par
\refstepcounter{exercise}\label{q:2:8}\par\noindent\textbf{习题\theexercise\quad 竞赛拓展} $f\in C^2[a,b]$有三个互异零点，证明$f\prime\prime$至少有一个零点。\quad\hyperref[a:2:8]{[答案]}\par
\refstepcounter{exercise}\label{q:2:9}\par\noindent\textbf{习题\theexercise\quad 竞赛拓展} 证明$\ln x\le x-1$，$x>0$。\quad\hyperref[a:2:9]{[答案]}\par
\section*{章末方法速查}\addcontentsline{toc}{section}{章末方法速查}
\par\noindent\begin{tabularx}{\textwidth}{p{23mm}XXX}\toprule 方法&适用条件&典型信号&易错点\\\midrule
链式法则&复合可导&内外函数&漏乘内导\\隐函数&$F_y\ne0$&隐式关系&漏掉$y\prime$\\Rolle&端点同值&导数零点&端点不连续\\Taylor&足够光滑&高阶接触&忽略余项\\
\bottomrule\end{tabularx}
